\documentclass[11pt,letterpaper]{article}

% ============================================================================
% PACKAGES
% ============================================================================
\usepackage[margin=1in, top=1.5in, headheight=85pt]{geometry}
\usepackage{graphicx}
\usepackage{fancyhdr}
\usepackage{lastpage}
\usepackage{enumitem}
\usepackage{xcolor}
\usepackage{hyperref}
\usepackage{tabularx}
\usepackage{listings}
\usepackage{courier}

% ============================================================================
% DOCUMENT VARIABLES
% ============================================================================
\newcommand{\UniqueID}{DM-2026-001}
\newcommand{\DocumentDate}{January 21, 2026}
\newcommand{\AuthorName}{SendCUIEmail Development Team}
\newcommand{\AuthorTitle}{Software Engineering}
\newcommand{\ToField}{Project Record}
\newcommand{\SubjectField}{Cross-Platform PowerShell Support Decision}
\newcommand{\OPRField}{SendCUIEmail Project}

% ============================================================================
% DOCUMENT CONFIGURATION
% ============================================================================
\hypersetup{
    colorlinks=false,
    pdfborder={0 0 0},
    pdftitle={DM-2026-001 Cross-Platform PowerShell Support},
    pdfauthor={SendCUIEmail Project},
}

\setlength{\parindent}{0pt}
\setlength{\parskip}{6pt}

\setlist[enumerate]{leftmargin=0.5in, labelwidth=0.25in, labelsep=0.1in, align=left, itemsep=12pt, topsep=12pt}
\setlist[enumerate,1]{label=\arabic*.}
\setlist[itemize]{leftmargin=0.5in, itemsep=3pt, topsep=6pt}

% Code listing style
\lstset{
    basicstyle=\small\ttfamily,
    breaklines=true,
    frame=single,
    backgroundcolor=\color{gray!10},
}

% ============================================================================
% HEADER AND FOOTER
% ============================================================================
\pagestyle{fancy}
\fancyhf{}
\fancyhead[L]{\textbf{SendCUIEmail}}
\fancyhead[R]{\parbox[b]{3.5in}{\raggedleft\large\textbf{Decision Memorandum}}}
\renewcommand{\headrulewidth}{0pt}
\fancyfoot[L]{\small \UniqueID}
\fancyfoot[C]{\small Page \thepage\ of \pageref{LastPage}}
\fancyfoot[R]{\small \DocumentDate}
\renewcommand{\footrulewidth}{0.4pt}

\newcommand{\dmsection}[1]{\item \textbf{#1}\par}

% ============================================================================
% BEGIN DOCUMENT
% ============================================================================
\begin{document}

\hfill \UniqueID

\vspace{0.25in}

\underline{\hspace{2.5in}}\\[2pt]
\AuthorName\\
\AuthorTitle

\vspace{0.25in}

\textbf{DATE:} \DocumentDate

\textbf{TO:} \ToField

\textbf{SUBJECT:} \SubjectField

\textbf{OFFICE OF PRIMARY RESPONSIBILITY:} \OPRField

\vspace{0.25in}

\begin{enumerate}

% ----------------------------------------------------------------------------
\dmsection{Purpose}

This memorandum documents the decision to implement cross-platform support for SendCUIEmail, enabling the encryption and decryption scripts to run on both Windows PowerShell 5.1+ and PowerShell Core 7+ (macOS and Linux).

% ----------------------------------------------------------------------------
\dmsection{Background}

SendCUIEmail was initially designed for Windows 11 environments to encrypt files for secure email transmission of CUI (Controlled Unclassified Information). The tool uses AES-256-CBC encryption with PBKDF2 key derivation, which are standard cryptographic algorithms available in .NET across all platforms.

During development, a requirement emerged to enable developers to test the core encryption/decryption logic on macOS development workstations without requiring a Windows virtual machine. This capability improves development velocity and enables continuous integration testing on non-Windows platforms.

% ----------------------------------------------------------------------------
\dmsection{Technical Analysis}

The primary compatibility issue between Windows PowerShell and PowerShell Core is the \texttt{SecureString} to plain text conversion. The standard Windows approach uses Marshal methods:

\begin{lstlisting}[language=PowerShell]
$bstr = [Runtime.InteropServices.Marshal]::SecureStringToBSTR($password)
$plain = [Runtime.InteropServices.Marshal]::PtrToStringAuto($bstr)
[Runtime.InteropServices.Marshal]::ZeroFreeBSTR($bstr)
\end{lstlisting}

This approach fails on non-Windows platforms because \texttt{PtrToStringAuto} behaves differently. The cross-platform alternative uses the \texttt{NetworkCredential} class:

\begin{lstlisting}[language=PowerShell]
$plain = [System.Net.NetworkCredential]::new('', $password).Password
\end{lstlisting}

\textbf{Platform Detection:} PowerShell Core exposes the \texttt{\$IsWindows}, \texttt{\$IsMacOS}, and \texttt{\$IsLinux} automatic variables. Windows PowerShell (Desktop edition) does not define these, so we detect the platform using:

\begin{lstlisting}[language=PowerShell]
$IsWindowsPlatform = $PSVersionTable.PSEdition -eq 'Desktop' -or $IsWindows
\end{lstlisting}

% ----------------------------------------------------------------------------
\dmsection{Scope of Changes}

The following modifications were made to support cross-platform operation:

\begin{itemize}
    \item Added \texttt{ConvertFrom-SecureStringPlain} helper function that uses platform-appropriate conversion
    \item Added \texttt{\$IsWindowsPlatform} detection variable
    \item Updated \texttt{Get-SecurePassword} in Encrypt.ps1 to use the helper function
    \item Updated \texttt{Get-Password} in Decrypt.ps1 to use the helper function
    \item Added platform check in \texttt{Generate-OutlookEmail} to skip .msg generation on non-Windows
    \item Updated script documentation to note cross-platform support
\end{itemize}

% ----------------------------------------------------------------------------
\dmsection{Platform-Specific Limitations}

The following features remain Windows-only:

\begin{itemize}
    \item \textbf{Outlook .msg Generation:} Requires Outlook COM automation, only available on Windows with Microsoft Outlook installed
    \item \textbf{Batch File Launchers:} The .bat files are Windows-specific; macOS/Linux users must invoke PowerShell directly
\end{itemize}

These limitations are acceptable because:
\begin{itemize}
    \item The target deployment environment is Windows 11
    \item macOS/Linux support is primarily for development and testing
    \item Core encryption/decryption functionality works identically across platforms
\end{itemize}

% ----------------------------------------------------------------------------
\dmsection{Security Considerations}

\textbf{No reduction in security posture:}
\begin{itemize}
    \item The cryptographic algorithms (AES-256-CBC, PBKDF2-HMAC-SHA256) are identical across platforms
    \item The encrypted file format is unchanged and interoperable
    \item Files encrypted on Windows can be decrypted on macOS/Linux and vice versa
    \item The \texttt{NetworkCredential} approach for SecureString conversion is a documented .NET pattern
\end{itemize}

% ----------------------------------------------------------------------------
\dmsection{Decision}

\textbf{APPROVED:} Implement cross-platform support as described above.

\textbf{Rationale:}
\begin{itemize}
    \item Enables faster development iteration without Windows VM dependency
    \item Supports CI/CD testing on multiple platforms
    \item No security impact---cryptographic operations are unchanged
    \item Graceful degradation for Windows-only features (Outlook .msg)
    \item Maintains full compatibility with existing encrypted files
\end{itemize}

% ----------------------------------------------------------------------------
\dmsection{Testing Requirements}

The following test scenarios must pass on both Windows and macOS/Linux:

\begin{itemize}
    \item Encrypt single file → Decrypt → Verify contents match
    \item Encrypt multiple files → Decrypt all → Verify contents match
    \item Wrong password rejection
    \item Empty file handling
    \item Binary files with null bytes
    \item Cross-platform interoperability (encrypt on macOS, decrypt on Windows)
\end{itemize}

\end{enumerate}

\end{document}
